\documentclass[11pt,a4paper]{ctexart}
\usepackage[margin=24mm]{geometry}
\usepackage{amsmath,amssymb,mathtools,bm,booktabs,array,microtype,hyperref}
\hypersetup{colorlinks=true,linkcolor=blue,urlcolor=blue}
\newcommand{\SO}{\mathrm{SO}}
\newcommand{\SE}{\mathrm{SE}}
\newcommand{\Diag}{\operatorname{Diag}}
\newcommand{\SceneBA}{\textsc{SceneBA}}
\title{\SceneBA：不确定性感知的离散--连续场景束调整\\
\large 数学建模与可实现推导}
\author{Lumenarium Research}
\date{2026年7月30日\quad v0.2}

\begin{document}
\maketitle
\begin{abstract}
\SceneBA 将单图三维场景重建从串行的``解析--检索--姿态--布局''流水线，
改写为离散场景假设与连续几何参数的联合后验推断。本文明确变量、因子图、
对称旋转、连续条件优化、Laplace 边缘证据、受限离散搜索，以及可证伪的
Top-$K$ oracle 实验。当前推导足以实现 oracle audit 和 bounded-beam 原型；
因子权重、候选阈值与统计非劣界仍须在独立开发集上标定。
\end{abstract}

\section{符号：$T_i,R_i,S_i$ 到底是什么？}
对第 $i$ 个物体，定义离散状态与连续状态
\begin{equation}
 z_i=(A_i,G_i,Q_i),\qquad
 x_i=(T_i,R_i,\ell_i),\qquad \ell_i=\log S_i .
\end{equation}
各符号含义如下。
\begin{itemize}
  \item $A_i$：资产身份，即 local/DeepSearch Top-$K$ 候选中的某个三维模型。
  \item $G_i=(P_i,\rho_i)$：场景图状态，其中 $P_i$ 是 parent/support，
  $\rho_i$ 是 \texttt{on}、\texttt{inside}、\texttt{wall}、\texttt{hang} 等关系。
  \item $Q_i$：离散姿态或对称模式，例如椅子的正反候选、四重旋转等价类。
  \item $T_i=(t_x,t_y,t_z)^\top\in\mathbb R^3$：物体局部原点在世界坐标系中的
  三维平移，单位通常为米。
  \item $R_i\in\SO(3)$：物体局部坐标系到世界坐标系的三维旋转矩阵，
  满足 $R_i^\top R_i=I,\det R_i=1$。
  \item $S_i=(s_x,s_y,s_z)^\top\in\mathbb R_{>0}^3$：沿三个局部轴的正尺度；
  等比例缩放是 $S_i=s_i\bm 1$。优化 $\ell_i=\log S_i$ 可自动保证尺度为正。
\end{itemize}
因此，物体到世界坐标的齐次变换为
\begin{equation}
 M_i(T_i,R_i,S_i)=
 \begin{bmatrix}
 R_i\Diag(S_i)&T_i\\
 \bm 0^\top&1
 \end{bmatrix}.
\end{equation}
对局部点 $v\in\mathbb R^3$，其世界坐标是
\begin{equation}
 v^{\mathrm w}=R_i\Diag(S_i)v+T_i.
\end{equation}
可选全局变量包括相机 $C$、深度尺度 $\alpha_d$、地面平面 $\pi_f$ 和墙面集合
$\Pi$，合记为 $\gamma=(C,\alpha_d,\pi_f,\Pi)$。

\section{对称感知旋转}
设资产 $A_i$ 的有限旋转对称群为 $H_{A_i}\subset\SO(3)$。物理姿态不是单个
$R_i$，而是商空间中的等价类
\begin{equation}
 [R_i]=\{R_i h:h\in H_{A_i}\}\in\SO(3)/H_{A_i}.
\end{equation}
对参考旋转 $\widehat R_i$，定义
\begin{equation}
 d_{\mathrm{sym}}(R_i,\widehat R_i;A_i)
 =\min_{h\in H_{A_i}}
 \left\|\operatorname{Log}\!\left(\widehat R_i^\top R_i h\right)^\vee
 \right\|_2 .
\end{equation}
其中 $\operatorname{Log}:\SO(3)\to\mathfrak{so}(3)$，$(\cdot)^\vee$ 将反对称矩阵
还原为旋转向量。连续更新使用李代数局部参数
\begin{equation}
 R_i(\delta\omega_i)=R_i^{(0)}
 \operatorname{Exp}([\delta\omega_i]_\times).
\end{equation}
当前系统可以先采用 yaw-only 的一维版本；完整 $\SO(3)$ 版本不改变离散建模。

\section{后验分布与因子图}
令观测
\begin{equation}
 y=(I,D,\mathcal M,\Phi)
\end{equation}
包含输入图像、深度、实例掩码与缓存的 foundation features。对
$z=(z_1,\ldots,z_N)$、$x=(x_1,\ldots,x_N)$，定义
\begin{equation}
 p(z,x,\gamma\mid y)
 =\frac{1}{Z(y)}p(z)\,p(x,\gamma\mid z)
 \prod_{f\in\mathcal F}\psi_f(z_f,x_f,\gamma;y).
\end{equation}
取负对数得到能量
\begin{equation}
 E(z,x,\gamma;y)
 =-\log p(z)-\log p(x,\gamma\mid z)
 +\sum_{f\in\mathcal F}E_f(z_f,x_f,\gamma;y),
 \quad E_f=-\log\psi_f .
\end{equation}
联合 MAP 为
\begin{equation}
 (z^\star,x^\star,\gamma^\star)
 =\arg\min_{z,x,\gamma}E(z,x,\gamma;y).
\end{equation}

\section{观测、先验和物理因子}
\subsection{图像与深度}
设 $\widehat M_i$ 是渲染掩码、$M_i$ 是观测掩码：
\begin{align}
 E_{\mathrm{sil},i}
 &=1-\operatorname{IoU}(\widehat M_i,M_i),\\
 E_{\mathrm{feat},i}
 &=\frac1{|\Omega_i|}\sum_{u\in\Omega_i}
 \rho\!\left(1-
 \cos(\widehat\Phi_i(u),\Phi_I(u))\right),\\
 E_{\mathrm{depth},i}
 &=\frac1{|\Omega_i|}\sum_{u\in\Omega_i}w_i(u)
 \rho\!\left(\log\widehat D_i(u)-\log(\alpha_dD(u))\right).
\end{align}
两物体投影重叠时，可加入遮挡顺序因子
\begin{equation}
 E_{\mathrm{occ},ij}
 =\frac1{|\Omega_{ij}|}\sum_{u\in\Omega_{ij}}
 \rho\!\left(s_{ij}(u)
 [\widehat D_i(u)-\widehat D_j(u)]\right),
\end{equation}
其中 $s_{ij}(u)\in\{-1,+1\}$ 编码观测到的前后关系。

\subsection{资产、图结构和初始化先验}
\begin{align}
 E_{\mathrm{ret},i}(A_i)
 &=-\log p_{\mathrm{ret}}(A_i\mid I,M_i,\text{label}_i),\\
 E_{\mathrm{graph},i}(G_i)
 &=-\log p(P_i,\rho_i\mid\text{S1 evidence}),\\
 E_{\mathrm{init},i}
 &=\frac12(x_i-\mu_i)^\top\Sigma_i^{-1}(x_i-\mu_i).
\end{align}
支撑图限制为有根有向森林：普通物体至多有一个结构 parent，且禁止有向环。
$\Sigma_i$ 大表示 S3 不确定、允许较大修正；$\Sigma_i$ 小则保护可靠初始化。

\subsection{物理与语义因子}
复用已经验证的 LayoutVLM 因子：
\begin{equation}
 E_{\mathrm{phys}}
 =\lambda_cE_{\mathrm{collision}}
 +\lambda_sE_{\mathrm{contact}}
 +\lambda_pE_{\mathrm{plane}}
 +\lambda_bE_{\mathrm{boundary}}
 +\lambda_mE_{\mathrm{semantic}}.
\end{equation}
语义项包含 gated containment、\texttt{align\_with}、\texttt{point\_towards}
和有向距离。contact、plane、containment 和 boundary 的硬约束可通过
投影算子 $\Pi_{\mathcal C}$ 保持可行：
\begin{equation}
 x^{k+1}
 =\Pi_{\mathcal C(z)}
 \left(x^k-\eta_k\nabla_xE(z,x^k,\gamma^k;y)\right).
\end{equation}

\subsection{总能量}
\begin{align}
 E={}&\sum_i\big[
 \lambda_{\mathrm{sil}}E_{\mathrm{sil},i}
 +\lambda_{\mathrm{feat}}E_{\mathrm{feat},i}
 +\lambda_dE_{\mathrm{depth},i}
 +\lambda_{\mathrm{ret}}E_{\mathrm{ret},i}\nonumber\\
 &\qquad+\lambda_{\mathrm{init}}E_{\mathrm{init},i}
 +\lambda_gE_{\mathrm{graph},i}\big]
 +\sum_{i<j}\lambda_oE_{\mathrm{occ},ij}
 +E_{\mathrm{phys}}.
\end{align}
训练或验证时必须单独报告各原始因子，不能把内部加权总能量当作唯一质量指标。

\section{为什么不能只按优化后的最小 loss 排名？}
对固定离散假设 $z$，令 $u=(x,\gamma)$ 且
\begin{equation}
 \widehat u_z=\arg\min_u E(z,u;y).
\end{equation}
我们真正需要比较的是边缘后验
\begin{equation}
 p(z\mid y)=\int p(z,u\mid y)\,du,
\end{equation}
而非单点密度。令
\begin{equation}
 H_z=\left.\nabla_u^2E(z,u;y)\right|_{u=\widehat u_z}.
\end{equation}
在局部极小值处二阶展开：
\begin{equation}
 E(z,u;y)\approx E(z,\widehat u_z;y)
 +\frac12(u-\widehat u_z)^\top H_z(u-\widehat u_z).
\end{equation}
高斯积分给出 Laplace 近似
\begin{align}
 p(z\mid y)
 &\approx \frac{1}{Z(y)}
 \exp[-E(z,\widehat u_z;y)]
 (2\pi)^{d_z/2}\det(H_z+\lambda I)^{-1/2},\\
 \log p(z\mid y)
 &\approx -E(z,\widehat u_z;y)
 -\frac12\log\det(H_z+\lambda I)
 +\frac{d_z}{2}\log(2\pi)+C.
\end{align}
如果 $E$ 的定义没有吸收离散先验，则须额外加 $\log p(z)$。为避免重复计数，
本文统一采用以下实现约定：
\begin{equation}
 E(z,u;y)=E_{\mathrm{likelihood}}+E_{\mathrm{continuous\ prior}},
\end{equation}
并将离散先验单独保留。于是排序分数为
\begin{equation}
 \boxed{
 \mathcal S(z)=
 -E(z,\widehat u_z;y)
 -\frac12\log\det(H_z+\lambda I)
 +\log p(z)
 +\frac{d_z}{2}\log(2\pi)}
\end{equation}
且越大越好；等价成本 $\mathcal C(z)=-\mathcal S(z)$ 越小越好。

曲率项是 Occam factor：低 loss 但极尖锐、脆弱的解具有更小的积分体积；
在较宽连续区域内都能解释观测的假设更可信。$\lambda>0$ 处理规范自由度和平坦模态。
不同维数的假设必须保留 $d_z\log(2\pi)/2$ 并使用一致先验测度。

\paragraph{近似计算。}
第一版无需构造稠密 Hessian。若残差向量为 $r(u)$、Jacobian 为 $J$，可用
\begin{equation}
 H_z\approx J^\top WJ+\lambda I
\end{equation}
的 Gauss--Newton 近似；再用稀疏 Cholesky、对角近似，或 stochastic Lanczos
估计 $\log\det H_z$。必须通过 ablation 比较 raw minimum、diagonal evidence
与更精确 evidence。

\section{受限的离散--连续推断}
完整枚举是组合爆炸，采用保留 incumbent 的不确定性门控 beam search：
\begin{enumerate}
  \item 从当前 Top-1 流水线结果初始化。
  \item 用 retrieval margin、parent entropy、symmetry ambiguity、渲染残差和
  物理残差计算风险。
  \item 选取最高风险物体及其 parent/children 构成局部 active subgraph。
  \item 提出 Top-3 assets、Top-2 parents 与少量 symmetry modes。
  \item 每个提案进行 30--50 步条件连续优化。
  \item 用 Laplace evidence 排名，保留 4--8 个 beam states。
  \item 冻结低风险子图，直到证据增益或算力预算耗尽。
  \item 对获胜状态运行 LayoutVLM400，或后续使用 $\{30,100,400\}$ router。
\end{enumerate}
保留 incumbent 可保证已评估 surrogate/evidence 不变差，但不保证 GT 指标单调提高。

\section{可辨识性、适用边界与实验闭环}
单视图中尺度和深度可互换，遮挡几何弱可观，对称物体存在多解。因此需要：
metric depth 或 floor/support priors、显式对称商空间、概率校准、初始化协方差，
以及对多模态后验的报告。

在实现完整求解器前，必须执行：
\begin{enumerate}
  \item asset Recall@$1/3/5/10$ 与 MRR；
  \item parent Recall@$1/2$；
  \item symmetry-aware pose oracle；
  \item asset-only、parent-only、symmetry-only 和 joint oracle；
  \item GT-free verifier 能恢复多少 oracle gap。
\end{enumerate}
建议 go/no-go 条件为 Recall@5 相比 Recall@1 至少提高 10 个百分点，或 joint oracle
使 rotation AUC 提高 0.05、translation AUC 提高 0.03、parent accuracy 提高 0.05
中的任意一项成立。最终求解器应恢复至少 40\% 的 oracle gap，展开不超过约
25--30\% 的物体，运行时间不超过 LayoutVLM400 的两倍。

\section{当前可以声称什么}
当前数学定义已经闭合并足以编码：变量域明确，后验和能量一致，对称性被商空间处理，
固定离散状态下的连续优化与离散边缘比较均有定义，搜索预算有界，并有可证伪实验。
但它还不是已被实验验证的最终算法。Top-$K$ 覆盖率、evidence 排名能力、因子权重和
SceneBA 对 Paper30 的真实增益仍需实验确定。

\end{document}
